%=== CHAPTER SIX ===
%% Rewritten 2026-09-18 against the deployed eventless interface.

\chapter{Conclusion and Future Work}
\label{ch:conclusion}
\begin{spacing}{1.5}
\setlength{\parskip}{0.22in}

\section{Summary}

This dissertation has treated the quadrotor payload problem as one coupled
estimation-and-control task with a specific structural difficulty at its centre.
The predictive model changes at grasp and at release in three ways: mass changes
the thrust-to-acceleration gain and the hover thrust, attachment geometry
increases roll and pitch inertia, and the horizontal centre-of-mass offset
creates a standing trim moment. Treating these as one undifferentiated
disturbance obscures that they act through different physical channels and are
identifiable to different degrees.

The structural difficulty is that the natural parameterisation of those changes
--- a payload mass together with a grasp offset --- is singular at release.
Chapter~\ref{ch:model} showed that as $\mP\to0$ the offset $\bm r_{xy}$ becomes
undefined, that holding it fixed leaves a phantom trim moment in the model, and
that the only way an optimiser can cancel that phantom is by corrupting the mass
estimate. The remedy adopted is to estimate the horizontal \emph{first mass
moment} $\svec=\mP\bm r_{xy}$ --- the product the equations of motion actually
contain --- alongside the composite mass. From that pair the centre-of-mass
offset follows as $\cvec=\svec/\mT$ with no division by the payload mass and no
knowledge of the grasp geometry, the off-diagonal inertia follows with the
payload mass cancelling exactly, and the dominant roll/pitch increment follows
from the composite mass and a weak vertical-arm prior whose error was measured to
be immaterial.

Two further reductions were necessary and were quantified rather than assumed.
The second-order horizontal inertia terms carry a factor $1/\mP$ in the new
coordinates and were discarded at a cost of \SI{2.6}{\percent} of the increment.
The remaining mass-dependent coefficient was shown to give the optimiser a lever
by which raising the mass estimate dilutes a phantom moment --- a mechanism that
drove the estimate \SI{53}{\percent} high in replay --- and was frozen within each
window to break that derivative while still tracking the true mass between
windows.

Around this model, Chapter~\ref{ch:framework} built a 16-state moving horizon
estimator driven by a wrench reconstructed from rotor speeds, a single atomic and
health-gated interface carrying that estimate to the controller, and a lifecycle
in which \emph{no attach or release notification is evidence}. The controller
issues the release command itself and does not consult it. The model is cleared
only by persistent estimator confidence; a release that cannot be confirmed
becomes \textsc{unresolved}, a hold in which the vehicle brakes smoothly to hover
and keeps waiting, rather than an assumption that the payload has gone.

In software-in-the-loop, Chapter~\ref{ch:evaluation} reported that the complete
interface confirmed release in 32 of 36 eligible flights through three distinct
confirmation paths, with median delays of \SI{3.46}{\second} in manoeuvre and
\SI{16.56}{\second} after braking to hover. Under a \SI{4}{\second} delayed
detachment, persistent confirmation avoided the early model clearing of a
command-triggered baseline in 5 of 5 pairs. The eventless interface additionally
detected 10 of 18 unsignaled payload losses that a command-armed confirmer cannot
see, at the cost of a detector exposed throughout loaded flight --- a cost that
produced no false declaration in 209 loaded flights.

\section{Principal Findings}

\begin{enumerate}
\item \textbf{The first mass moment, not the offset, is the estimable payload
geometry.} The equations of motion contain only the product. Estimating the
factors separately adds an unidentifiable direction and a singularity at exactly
the operating point of interest.

\item \textbf{Payload \emph{presence} must not be encoded in the mass channel.}
When it is, the only way for the estimator to express ``the load has gone'' is to
move the mass, which is the one quantity the controller most needs to be correct.
Augmenting $\svec$ gives presence its own channel.

\item \textbf{A parameter that appears in a denominator will be abused by an
optimiser.} Both pathologies found in this work --- the $1/\mP$ coefficient of
the discarded inertia terms, and the mass in the denominator of the ghost-moment
expression --- are instances of the same phenomenon. The fix in both cases was
structural, not a matter of tuning.

\item \textbf{An adaptive loop must be driven by a signal that is not a product
of the adaptation.} Feeding the controller's intended thrust to the estimator
makes the mass unobservable by construction; the rotor-speed reconstruction
breaks the circularity, and extending it to the yaw-reaction channel removed the
last place where the estimator could learn the controller's intent.

\item \textbf{Release evidence must come from the rotational channel.} A
collective-thrust residual cannot distinguish a payload change from a trajectory
change or a manoeuvre; a rule that omitted mandatory moment evidence released
prematurely in 3 of 9 flights, up to \SI{51}{\second} early. Making moment
collapse mandatory costs misses, which are handled as \textsc{unresolved} holds
rather than by lowering the evidence bar.

\item \textbf{A ratio needs a reference, and the reference must be formed
autonomously and then frozen.} A running maximum absorbs the very estimation
errors the detector exists to reject.

\item \textbf{Timeout is not proof.} Treating an unconfirmed release as complete
reintroduces the failure the design exists to prevent. \textsc{unresolved} is the
state that makes eventlessness deployable.

\item \textbf{Payload adaptation must update the objective as well as the state
equation.} A stale equilibrium-input baseline produces a steady-state error that
gets \emph{worse} as the input weight is increased --- a signature that
identifies the fault and that no weight tuning can remove.

\item \textbf{The interface is severely sensitive to thrust-map calibration.} A
\SI{\pm5}{\percent} error in the estimator's thrust coefficient defeated payload
recognition in both directions. This is the principal barrier to hardware
transfer, and it is a direct consequence of recovering an absolute mass from an
absolute thrust.
\end{enumerate}

\section{Limitations}

All system-level evidence is software-in-the-loop, and the simulator's plant
generates thrust with the same square law the estimator inverts. Rotor speeds are
exact; speed noise, transport delay and induced inflow are untested.

The attachment is rigid. Gripper compliance, pendulum modes, flexible and
sloshing loads are not modelled. A sustained vertical gust remains intrinsically
confounded with a payload change on a collective-thrust residual, and wind is not
evaluated.

Within the interface itself, the dominant failure is not a wrong decision but an
unavailable one: in 13 of 32 completed flights the frozen moment reference never
formed, so the ratio test never ran, and those flights are also the population of
seven of the eight unsignaled-loss misses. Out-of-envelope first-moment samples
are excluded from votes but still reach the controller's centre-of-mass parameter
through filtering. The release command is not gated on estimator health, and one
divergence is directly attributable to that omission.

The heaviest condition, \SI{0.30}{\kg} at \SI{4}{\meter\per\second}, sits at the
vehicle's torque limit rather than at an operating point: pitch torque saturated
within \SI{98}{\percent} of diagnostic windows there. Safe release
\emph{semantics} are claimed at that limit; reliable flight is not.

Finally, the baseline set is thin. The reported comparisons isolate release
semantics only. No comparison against a fixed-parameter controller, a
lumped-disturbance adaptive controller, or a formal change-detection procedure
under a matched false-alarm constraint is included, although implementations of
the first two exist in the codebase.

\section{Future Work}

\paragraph{Reference formation.} The single highest-value improvement is to make
the loaded moment reference form reliably. Roughly \SI{40}{\percent} of flights
never produced one, and that single deficiency accounts for the brake-to-hover
delay mode, most unsignaled-loss misses, and the confidence-only confirmations.
Candidates include forming the reference during the lift transient rather than
requiring a later quiet window, admitting a direction-only reference when the
magnitude is unstable, and exiting the manoeuvre early on the envelope-exceedance
counter that already correlates with the outcome at \SI{4}{\meter\per\second}.

\paragraph{Calibration robustness.} Because the mass channel consumes an absolute
thrust, its accuracy is bounded by the accuracy of $k_f$. Two directions are
open: propagating a calibration uncertainty into the estimator weights so that
the interface degrades gracefully rather than at a threshold, and exploiting the
\emph{calibration-invariant} structure already present in the centre-of-mass
inversion, in which $k_f$ cancels, as an independent cross-check on the mass
channel.

\paragraph{Health-gated release.} The release command should not be issued while
the estimator is unhealthy. This is a small change with a demonstrated failure
behind it.

\paragraph{Baselines and detection theory.} A matched-false-alarm comparison
against CUSUM and related procedures on the same signals would establish whether
the hand-built persistence rules are near the achievable operating curve or far
from it. A comparison against a fixed-parameter and an $\mathcal{L}_1$-adaptive
controller would separate the benefit of the payload model from the benefit of
adaptation in general.

\paragraph{Modelling extensions.} The yaw inertia increment and the second-order
horizontal terms could be restored if a formulation is found that avoids the
$1/\mP$ degeneracy --- for example by regularising the payload mass away from
zero only inside those terms. A random-walk treatment of the payload states would
remove the remaining one-window transient, provided the Hessian scaling problem
it caused can be solved.

\paragraph{Hardware transfer.} Transfer should proceed in stages: motor-stand
calibration of $k_f$ and $k_m$ against measured thrust and reaction torque;
tethered hover with the estimator running open-loop; centred payload pickup;
eccentric pickup; and only then dynamic flight. Simulated rotor speeds must be
replaced by electronic-speed-controller telemetry, and the noise, delay and
quantisation of that telemetry re-measured rather than assumed. A supervisor
monitoring solver status, estimator health and input margin should gate every
stage.

\section{Closing Remark}

The lesson that generalises beyond this vehicle is about the relationship between
an estimate and a decision. Online estimation is useful only when the estimated
quantity enters the controller in a physically and numerically consistent way ---
and a system that acts on events it has been \emph{told} about, rather than on
what it can \emph{observe}, has no way to notice when the two disagree. Choosing
a parameterisation that stays well posed through the transition, and refusing to
let a command stand in for evidence, are the two decisions that made the rest of
this architecture work.

\end{spacing}
\newpage
